\section{决策树把预测写成递归分区}
\label{sec:13-Machine-Learning/03-Trees-Ensembles-and-Neighbors/01}

你手头有一份贷款违约数据，年龄、收入、负债比、信用历史各有几列。线性回归跑完，每个系数讲的是一个全局故事：收入每高 1 单位，违约概率下降固定的量。但你盯着残差图看——年轻且高负债的客户被系统性低估，年长且低负债的被高估。你在模型中加交互项：年龄\(\times\)负债比、年龄\(\times\)收入……特征开始爆炸，每个交互项的解释也变得模糊。

\textbf{困境是}：很多真实规律并非全局一致，而是``在这个条件下适用这条规则，在另一个条件下适用另一条''。线性模型虽然可以通过显式构造交互项来适配，但成本是参数数量随组合数指数增长。

决策树走了一条完全不同的路径：不写一个覆盖全空间的公式，而是反复把空间一分为二，在每一片内部给一个简单常数预测。

\subsection{如果一个区域已经分好了，常数预测就是均值}

假设经过一系列切分，样本落到了某个叶节点，包含索引集合 \(R\)。该叶用常数 \(c\) 做预测，平方损失为

\[
\sum_{i\in R}(y_i-c)^2.
\]

求导置零立刻得到

\[
c_R = \frac1{|R|}\sum_{i\in R}y_i.
\]

叶节点均值是局部平方损失的最优常数。如果换成绝对损失，最优常数变成叶内中位数——切分规则不变，预测语义已经不同。损失函数的选择在树模型中同时影响切分标准和叶预测值，这两层耦合不能独立设定。

分类树的逻辑类似，只是把均值换成类别比例：

\[
\widehat p_{Rk} = \frac1{|R|}\sum_{i\in R}\ind\{y_i=k\}.
\]

最终类别可以取最大比例的那一个，也可以乘以错误分类的成本矩阵另作选择。需要留心的是，小叶节点上的比例是训练样本的频率，可能高度离散——比如在一个只有 3 个样本的叶里，比例只可能是 \(0, \frac13, \frac23, 1\) 中的某一个。它不是一个校准好的概率估计。

\subsection{不纯度：衡量一次切分清理了多少混合}

既然要递归切分，就需要一个标准来判断哪个切分更好。给定当前节点的类别比例 \(p_1,\ldots,p_K\)，什么叫``不纯''？

两个最常用的答案是 entropy：

\[
H(p)=-\sum_{k=1}^{K}p_k\log p_k
\]

和 Gini impurity：

\[
G(p)=1-\sum_{k=1}^{K}p_k^2.
\]

两者在节点中只有一个类别时为零（完全纯净），在各类均匀时达到最大。候选切分把父节点 \(R\) 分成左子 \(R_L\) 和右子 \(R_R\)，不纯度下降为

\[
\Delta I = I(R) - \frac{|R_L|}{|R|}I(R_L) - \frac{|R_R|}{|R|}I(R_R).
\]

子节点的不纯度按样本比例加权——防的是你通过切出一个只包含 1 个纯样本的极小节点来伪造巨大的``提升''。回归树把 \(I(R)\) 替换为叶内平方误差，逻辑完全相同：加权方差减少。

entropy 和信息论中的条件熵形式上一致，但这里有一个容易踩的坑：训练集上的信息增益会对高基数特征产生系统性偏好。比如把客户 ID 当成特征，每个 ID 只出现一次——按 ID 切分可以把每个训练样本单独成叶，不纯度降到零。但这对新客户没有可复用的结构。不纯度下降衡量的是训练数据上的经验分区质量，不是泛化能力的保证。

\subsection{贪心生长永远不会是最优树，但可以是有用的树}

寻找给定深度下的全局最优树是 NP-hard。各算法退而求其次：在每个节点，枚举所有特征和所有可行阈值，计算不纯度下降，选出当前最佳的那一刀，然后递归处理两个子节点。

这意味着早期的切分决定了后续可以表达哪些分区模式——第一次切了``年龄<30''之后，左子树只能在年轻人群中继续细分，而无法再和老年人群混合做更精细的划分。当前最优不等于全局最优。

由此也解释了决策树对数据扰动的剧烈敏感性。两个候选切分的收益非常接近时，少数样本的增减或数值微小变化就可能交换它们的排名。一旦根节点的选择变了，整棵子树跟着大换血。预测函数是分段常数的，但学习算法远不是输入数据的平滑函数——这也是后面集成方法要专门处理的问题。

\subsection{深度控制的本质是局部样本量}

你当然可以不断往下分裂，直到每个训练样本独占一叶，训练误差降到零。代价是叶内估计只用到了极少样本，方差极大——下一次采样，同一个叶可能因为一两个点就变成完全不同的均值。

四种常见的约束手段分别在限制不同的东西：

\begin{itemize}
\tightlist
\item
  \textbf{最大深度}：限制从根到叶的最长路径长度；
\item
  \textbf{最小叶样本数}：规定一个叶至少有多少个样本才能被创建；
\item
  \textbf{最小不纯度下降}：切分必须带来足够的信息增益，否则不切；
\item
  \textbf{剪枝}：先长成一棵大树，再自下而上地剪掉贡献不足的子树。
\end{itemize}

代价复杂度剪枝的形式最清晰：

\[
\widehat R(T) + \alpha|\mathcal L(T)|,
\]

其中 \(T\) 是树，\(\mathcal L(T)\) 是叶集合。第一项奖励拟合，第二项为每个叶节点收取复杂度税。参数 \(\alpha\) 控制你愿意为多一个叶付出多少训练误差的减少。

但 \(\alpha\) 不能从训练集上选——你仍然需要一块没参与拟合的数据来决定剪多少。先看测试集挑深度，再报告同一份测试集的分数，等于让测试集做了超参数选择，重新引入选择乐观性。保留一份独立的验证集，或者用交叉验证选 \(\alpha\)，只对最终树报告测试结果。

\subsection{轴对齐分区是一把双刃剑}

普通决策树的每次切分检查的是单个特征与一个阈值的比较：\(x_j \le t\)。这意味着所有分区边界都平行于坐标轴——矩形区域。

好处很明显：它能自然捕获阈值效应（``收入过了一个门槛就完全不同了''），不需要对原始特征做标准化，不同量纲的变量可以平等竞争切分位置。交互效应也不需要显式构造——如果某个子区域里年龄和负债比产生了交互，连续两个切分就能逼近它。

限制同样明显：一条 \(y = 2x + 1\) 的对角线，树只能用许多小矩形拼出一个楼梯来近似。旋转一下坐标系，同一份数据可能需要完全不同的树结构。也别忘了，树虽然不需要为欧氏距离做尺度统一，但类别编码的方式、缺失值的填补策略、时间上的数据泄漏和高基数特征的存在仍会大幅影响切分的选择。预测外推时，叶节点没有参数模型那种趋势延续的能力——它把训练数据边界内的样本做平均，边界之外就返回最邻近叶的常数。

一棵树的透明在于可以逐路径读出决策逻辑，但这不代表分区具有因果关系，也不代表一棵很深的树整体容易理解。后续集成章节会降低单树的方差，代价是放弃了这种逐路径的可读性——你在效率和可解释之间做了一次结构性的权衡。
